\section{The substitution construction and its dynamical properties} 
\label{sec:comp}

%
%
%
%
%

\subsection{Pre-substitution}

Let $\mathbf a=(a_i)_{i\geqslant 0}=a_0,a_1,a_2,\ldots$ be a sequence of nonnegative integers that satisfies the following additional properties.

\begin{enumerate}[label={(A\arabic*)}]
\item \label{def:a1} The sequence $\mathbf a$ is bounded, where we set $N:=\max a_i$.
\item \label{def:a2} $a_0\neq 0$.
\item \label{def:a3} The runs of zeros in $\mathbf a$ are bounded. That is, there exists an integer $C>0$ such that if $a_i=a_{i+1}=\ldots=a_{i+k}=0$, then $k<C$.
\end{enumerate}

%


\begin{definition} \label{def:substitution}


Let $\mathbf a\in \left\{0,1,\ldots,N\right\}^{\mathbb{N}_0}$ be given and let $[i]\in\mathbb{N}_0$ be a letter in our pre-alphabet. We define the {\it pre-substitution} $\varrho^{ }_\mathbf a\colon \N_0\to \N_0^{+}$ via
$$
\begin{array}{l}
\varrho^{ }_\mathbf a([0])=[0]^{a_0}[1], \text{ and}\\
\varrho^{ }_\mathbf a([i])=[0]^{a_i}[i-1][i+1] \text{ for }i>0.
\end{array}
$$
\end{definition}

Here $[0]^{a_i}$ denotes $a_i$ zeros in a row.
We call the rule $\varrho^{ }_{\mathbf{a}}$ a pre-substitution because we will later extend it to a full substitution on a compact alphabet. This is possible if one embeds $\mathbb{N}_0$ properly. 
To this pre-substitution, one can associate the corresponding infinite matrix 
$$\mathbf A=\left( 
\begin{array}{cccccc}
a_0& a_1+1& a_2 & a_3 & a_4 & \dots\\
1 & 0 & 1 & 0 & 0 & \ldots\\
0 & 1 & 0 & 1 & 0 & \ldots\\
0 & 0 & 1 & 0 & 1 & \ldots\\
0 & 0 & 0 & 1 & 0 & \ldots\\
\vdots & \vdots & \vdots & \vdots & \vdots & \ddots
\end{array}
\right),$$
whose entries are defined by $\mathbf{A}_{ij}$ being the number of $[i]$ in $\varrho_{\mathbf{a}}([j])$. This matrix will become relevant later when we compute for statistical and geometric properties associated to $\varrho$. 
This construction can be seen as a generalization of \cite[Sec~4.4]{MRW}. This example is a specific case of such a substitution where one has $a_i=1$ for all $i\geq 0$.



\subsection{Compactification and its properties}%

In order to use the results of \cite{MRW2}, the alphabet has to be compact in some topological space. Here, we use the boundedness of the entries of $\mathbf{a}$ to define an appropriate compactification of $\mathbb{N}_0$. 


Consider the full shift $\mathcal S=\{*,0,1,2,\ldots,N\}^\Z$ over $N+2$ letters. This is a compact topological space under the local topology.  More specifically, for every $i\geq 0$, the $i$th non-trivial neighborhood (also called a cylinder set) of $\mathbf b$ in $\mathcal S$ consists of those elements of $\mathcal S$ that have the same subsequence $b_{-i}\ldots b_{-1}\fbox{$b_0$}b_1\ldots b_{i}$ around the origin. Here and in the sequel, the boxed term will always specify the location of the origin. 

%

It is well known that $\mathcal S$ is a metric space where the distance $d_\mathcal S(\mathbf b,\mathbf b')$ between $\mathbf b,\mathbf b'\in \mathcal S$ is given by $\frac{1}{2^{i}}$, where $i$ is the smallest number in  $\N_0$ such that the $i$th nontrivial neighborhoods of $\mathbf b$ and $\mathbf b'$ are different. In particular, if $\mathbf b$ and $\mathbf b'$ have different letters at the origin then $i=0$ and $d_\mathcal S(\mathbf b,\mathbf b')=1$.

%
%
%
%


Fix a sequence $\mathbf a\in \left\{0,1,\ldots,N\right\}^{\mathbb{N}_0}$. We now define an embedding $\iota_{\ab}\colon \N_0\to \mathcal{S}$ depending on $\ab$, which we do first by defining $\mathbf 0:=\iota_{\ab}([0])$. The image of $[0]$ is defined pointwise via the rule
\[
\iota_{\ab}([0])_m:=\begin{cases}
a_m,& \text{ for } m\geqslant 0, \\
*,& \text{otherwise}.
\end{cases}
\]
As a bi-infinite sequence, one has 
$$\mathbf 0:=\ldots ***\fbox{$a_0$}a_1a_2a_3\ldots\,.$$  For every $i>0$, we define $\iota_{\ab}([i]):=\sigma^{i}(\mathbf 0)$, i.e., 
$$\mathbf i:=\ldots ***a_0a_1\ldots a_{i-1}\fbox{$a_i$}a_{i+1}a_{i+2}\ldots\,.$$
Here, $\sigma$ is the left shift map as before. 

The closure $\mathcal{A}:=\overline{\iota_{\ab}(\N_0)} \subset \mathcal S$ is the alphabet that we use in our further construction. For an element $\mathbf b\in \mathcal{A}\setminus \iota_{\ab}(\N_0)$, we will denote the corresponding letter of the alphabet as $[\infty_\mathbf b]$. Note that we will always obtain at least one $[\infty_{\mathbf b}]$ from the compactification. Thus the total {\it compact alphabet} $\mathcal{A}$ can be seen as the union of all isolated letters $\iota_{\ab}([i])$ and all relevant infinities $[\infty_\mathbf b]$. For brevity, for every $i\in \N_0$, we identify the image $\iota_{\ab}([i])$ in the alphabet with $[i]$ in the pre-alphabet and use the notation $[i]$ for the isolated letters, so 
$$\mathcal{A}=\{[i]\mid i\in \N_0\}\cup \{[\infty_\mathbf b]\mid\mathbf b\in \mathcal{A}\setminus \iota_{\ab}(\N_0)\}.$$

Before defining the substitution on the letters $[\infty_\mathbf b]$, we prove two properties of the corresponding sequences in $\mathcal S$.

\begin{lem}\label{lem:nostar}
Let $\mathbf{b}\in \mathcal{A}$. 
Then $\mathbf b\in \mathcal{A}\setminus \iota_{\ab}(\N_0)$ if and only if no entry of $\mathbf b$ is $*$.
\end{lem}
\begin{proof}
This follows directly from the definition of $\iota_{\ab}(\N_0)$ as $\iota_{\ab}(\N_0)=\left\{\sigma^{n}(\mathbf{0}),n\in\mathbb{N}_0\right\}$. For every $n$, only the entries to the left of the $(-n)$th position of $\sigma^n(\mathbf 0)$ are $*$s. This means that accumulation points of $\iota_{\ab}(\N_0)$ are not in the $\sigma$-orbit of $\mathbf{0}$  and must be bi-infinite sequences over $\left\{0,1,\ldots,N \right\}$.
\end{proof}

\begin{lem}\label{lem:shift}
The shift map is invertible on the accumulation points. 
That is, if one has $\mathbf b\in \mathcal{A}\setminus \iota_{\ab}(\N_0)$, then $\sigma(\mathbf{b}),\sigma^{-1}(\mathbf{b})\in \mathcal{A}\setminus \iota_{\ab}(\N_0)$. 
\end{lem}

\begin{proof}
Since $\mathbf b$ is an accumulation point, there exists an increasing sequence $(i_n)_{n\geqslant 0}$ of natural numbers such that the sequence $(\mathbf{i}_n)_n$ with ${\bf i}_n:=\sigma^{i_n}({\bf 0})$ converges to $\bf{b}$ in the local topology. Without loss of generality, assume that $i_n>1$ for all $n$. Note that the sequence $\left(\sigma^{i_n-1}(\bf{0})\right)_{n\geqslant 0}$ converges to $\sigma^{-1}(\bf{b})$ and all elements of the sequence are in $\iota_{\ab}(\N_0)$ because they are shifts of $\bf{0}$, which implies $\sigma^{-1}(\bf{b})\in \overline{\iota_{\ab}(\N_0)}$. The claim for $\sigma({\bf{b}})$ can be proved using similar arguments. 
\end{proof}


We then extend the substitution in Definition~\ref{def:substitution} to the image of  all $[\infty_\mathbf b]$ under $\varrho^{ }_\mathbf a$. 
For $$\mathbf b=\ldots b_{-2}b_{-1}\fbox{$b_0$}b_1b_2\ldots \in \mathcal{A}\setminus \iota_{\ab}(\N_0),$$ 
we define
\begin{equation}\label{eq:subs-extend}
\varrho^{ }_\mathbf a([\infty_\mathbf b])=[0]^{b_0}[\infty^{ }_{\sigma^{-1}(\mathbf b)}][\infty^{ }_{\sigma(\mathbf b)}].
\end{equation}
%

Lemmas \ref{lem:nostar} and \ref{lem:shift} ensure that this substitution is well defined as $b_0$ must be a number from $\{0, 1,\ldots,N\}$ and not $*$, and that the letters $[\infty^{ }_{\sigma(\mathbf b)}]$ and $[\infty^{ }_{\sigma^{-1}(\mathbf b)}]$ exist in $\mathcal{A}$.


%

%
\begin{prop}
The map $\varrho^{ }_\mathbf a\colon \mathcal{A}\to \mathcal{A}^{+}$ defined in Def.~\ref{def:substitution} and Eq.~\eqref{eq:subs-extend} is continuous and hence a substitution. 
\end{prop}
\begin{proof}
 This follows immediately from the construction. 
In particular, if $d_\mathcal S(\sigma^{i}(\mathbf{0}),{\bf{b}})\leq \frac{1}{2^m}$ for some $m\geq 1$ and $i>0$, then $\mathbf i$ and $\mathbf b$ coincide at least within $m$ letters from the origin. Thus
$d_\mathcal S(\sigma^{i-1}(\mathbf{0}),\sigma^{-1}(\mathbf{b}))\leq\frac{1}{2^{m-1}}$. Similarly,  $d_\mathcal S(\sigma^{i+1}(\mathbf{0}),\sigma(\mathbf{b}))\leq \frac{1}{2^{m-1}}$. 
Additionally, the condition on $\sigma^{i}(\mathbf{0})$ implies $a_i=b_0$. 
This means that the words $\varrho^{ }_{\mathbf{a}}([i])=[0]^{a_i}[i-1][i+1]$ and  $\varrho^{ }_{\mathbf{a}}([\infty_{\mathbf{b}}])=[0]^{b_0}[\infty_{\sigma^{-1}(\mathbf{b})}][\infty_{\sigma(\mathbf{b})}]$ both belong to $\mathcal{A}^{b_0+2}$ and are close in $\mathcal{A}^{+}$ in the disjoint union topology.

The cases of other letters from $\mathcal A$ are similar.
\end{proof}
%

\begin{thm}\label{thm:primrecquasi}
Let $\bf{a}$ be a sequence of nonnegative integers satisfying conditions \textnormal{\ref{def:a1}--\ref{def:a3}} and $\varrho^{ }_{\bf{a}}$ be the substitution on  $\mathcal{A}$ corresponding to $\bf{a}$.
The substitution $\varrho^{ }_\mathbf a$ satisfies the following.
\begin{enumerate}[label={{\rm (\roman*)}}]
\item $\varrho^{ }_\mathbf a$ is primitive.
\item $\varrho^{ }_\mathbf a$ is recognizable. 
\item The associated (normalized) substitution operator is quasicompact. 
\end{enumerate}
\end{thm}
\begin{proof}


We first prove that every letter $a\in \mathcal{A}$ contains $[0]$ in its level-$(C+1)$ supertile $\varrho_\mathbf a^{C+1}(a)$, where $C$ is the constant from \ref{def:a3}. For this purpose, note that $\varrho^{ }_\mathbf a([0])$ contains $[0]$ because $a_0\neq 0$ due to \ref{def:a2}. Thus, if a supertile contains $[0]$, then every further application of $\varrho^{ }_\mathbf a$ to that supertile will contain $[0]$ as well.

If $a\in \iota_{\ab}(\N_0)$, then $a=[i]$ for some $i$. Among the integers $a_i,a_{i+1},\ldots,a_{i+C}$, at least one is positive due to condition \ref{def:a3}. Suppose $a_{i+k}>0$, where $0\leq k \leq C$. Then $\varrho_\mathbf a^{k}([i])$ contains the letter $[i+k]$. Since $\varrho_\mathbf a^{k+1}([i])$ contains $[0]$, so does $\varrho_\mathbf a^{C+1}([i])$.

If $a\in \mathcal{A}\setminus \iota_{\ab}(\N_0)$, then $a=[\infty_\mathbf b]$ for an appropriate $\mathbf b=\ldots b_{-2}b_{-1}\fbox{$b_0$}b_1b_2\ldots$. Similarly, among the numbers $b_0,b_1,\ldots,b_C$ at least one is positive because $\mathbf b$ is a limit of some subsequence of $(\sigma^i(\mathbf 0))_i$. After that, the arguments are similar to the previous case as $\varrho_\mathbf a^{k}([\infty_\mathbf b])$ contains $[\infty_{\sigma^k(\mathbf b)}]$.


%

We proceed with the proofs for the three claims in the theorem.

For primitivity, we need to show that, for every non-empty open set $U$, there exists $k(U)\in\mathbb{N}$ such that for every letter $a\in \mathcal{A}$, $\varrho_{\bf{a}}^{k(U)}(a)$ contains a letter in $U$. 
We consider two types of open sets: the singletons consisting of isolated points $[j],j\in\N_0$, and balls $B_{\varepsilon}([\infty_{\bf{b}}])$ for some accumulation point $\bf{b}$. 

If $[j]$ is an isolated point, the claim  is straightforward for $U=\left\{[j]\right\}$ since $[j]\triangleleft\varrho^{j}_{\bf{a}}([0])$ and hence $[j]\triangleleft\varrho^{j+C+1}_{\bf{a}}(a)$ for any $a\in \mathcal{A}$.

Now let $\bf{b}$ be an accumulation point and $U=B_{\varepsilon}([\infty_{\bf{b}}])$ for some $\varepsilon>0$. Following the proof of Lemma~\ref{lem:shift}, there exists a sequence of isolated points $\left\{{\bf{i}}_{n}\right\}_{n\geqslant 0}$ such that for every $\varepsilon>0$, there is a $T$ such that $d_\mathcal S({\bf{i}}_n,\bf{b})<\varepsilon$ for all $n>T$. For the embedding, this means $[i_n]\in U$. Same as above, we pick any $n>T$ and then for every $a\in\mathcal{A}$, we get that $[i_n]\in\varrho_{\bf{a}}^{i_n+C+1}(a)$, which proves the claim.


Next, we show that the substitution $\varrho^{ }_{\bf{a}}$ is recognizable. Here, we use three facts: 
\begin{enumerate}
        \item Zeros are always at the start of a supertile. (However, there may be level-$k$ supertiles without zeros for some $k\geq 1$.)
        \item  The supertiles $\varrho_{\bf{a}}([0])$ and $\varrho_{\bf{a}}([1])$ contain exactly one non-zero letter each.
        \item  Any supertile other than $\varrho_{\bf{a}}([0])$ and $\varrho_{\bf{a}}([1])$  contains exactly two non-zero letters.
\end{enumerate}
  In particular, all three facts together imply 
 \begin{enumerate}
     \item Runs of consecutive zeros belong to one single level-1 supertile.
     \item A supertile without zeros has exactly two letters.
 \end{enumerate}
  Thus, we can proceed as follows in order to determine level-1 supertiles uniquely:
whenever we see a (maximal) block of zeros, then a level-1 supertile starts with that block and the whole block belongs to that one supertile. 
%

Given $x\in X_{\varrho^{ }_{\bf a}}$, we show that there is a unique way to decompose it as a  concatenation of level-1 supertiles. First, we look to the right of the origin and look for the minimal $n\in\mathbb{N}$ such that $x_{n-1}\neq [0]$ and $x_n=[0]$. This means $x_n$ is a beginning of a level-1 supertile. We iterate the process to the right and to the left of $x_n$, placing a cut each time we are in such a situation. Between each two consecutive cuts we have a block of zeros followed by some number of non-zero letters. Since every level-1 supertile without initial zeros has exactly two non-zero letters, the initial block of zeros form a level-1 supertile with one or two non-zero letters (depending on the parity of the number of non-zero letters between two consecutive cuts). From this, one gets the unique decomposition into level-1 supertiles. 

The process for level-$k$ decomposition for $k>1$ is similar, where now one has to look at positions where $x_{[n,n+|\varrho^{k-1}_{\bf{a}}([0])|-1]}=\varrho^{k-1}_{\bf{a}}([0])$ and $x_{[n-|\varrho^{k-1}_{\bf{a}}([0])|,n-1]}\neq\varrho^{k-1}_{\bf{a}}([0])$. Since the substitution is injective, the subwords between every two consecutive positions that have this property can be unambiguously cut into level-$k$ supertiles. 

Lastly, we show that the associated (normalized) substitution operator is quasicompact. Here, we want to show that $\varrho^{ }_{\bf{a}}$ satisfies the conditions of Proposition~\ref{prop:quasicompact}. We show that this is satisfied for $k=C+2$, where $C$ is the constant from property \ref{def:a3} and for $F=\{[0],[1],\ldots,[C+1]\}$. 

We look at the images of letters in $\mathcal{A}$ under $\varrho^{C+2}_{\bf{a}}$. Every application of $\varrho^{ }_{\bf{a}}$ at least doubles the number of letters and therefore for all $a\in\mathcal{A}$, $|\varrho^{C+2}_{\bf{a}}(a)|\geq 2^{C+2}$. However, if we apply $\varrho^{ }_{\bf{a}}$ that many times, then at least once one letter will be substituted with three or more letters. Indeed, if $a=[j]$ for some $j\geq 0$, then after first $C+1$ applications of $\varrho_{\bf{a}}$, we will see the letters $[j+1], [j+2], \ldots, [j+C+1]$ and at least one of $a_{j+1}, a_{j+2}, \ldots, a_{j+C+1}$ must be positive by property \ref{def:a3}. The corresponding letter is substituted by at least three letters. The case for $a=[\infty_\mathbf b]$ is similar. This implies $|\varrho^{C+2}_{\bf{a}}(a)|\geq 2^{C+2}+1$ and by Lemma \ref{lem:spec-rad}, $r(M)^{C+2}\geq 2^{C+2}+1$.

On the other hand, we note that for every $n\leq C+1$, the supertile $\varrho^{n}_{\bf{a}}([0])$ only contains letters from $F$, and the supertile $\varrho^{C+2}_{\bf{a}}([0])$ contains only one letter not from $F$. For every other letter $a\neq [0]$, we will repeatedly apply $\varrho^{ }_{\bf{a}}$ $C+2$ times to $a$ but we will color some letters with red in the process. The goal is to show that red letters cannot lead to letters outside of $F$ and there are not too many non-red letters.

Here are the rules for how we color letters with red inductively. The initial letter $a$ is not colored with red. If one applies $\rho^{ }_{\bf{a}}$ to a non-red letter $[i]$ (or appropriate infinity), then the initial zeros of $\rho^{ }_{\bf{a}}([i])$, if any, are colored with red. Additionally, if we apply $\rho^{ }_{\bf{a}}$ to a red letter, then all resulting letters are red. For the sake of lucidity we illustrate the coloring process by an example in Table \ref{tab:ex-red}.

\begin{table}
\[ \mathbf{a}=1,0,1,0,1,0,1,0,1,0,1,0,\ldots\]
\[ \begin{array}{cll}
0 & \mapsto & 01\\
1 & \mapsto & 02\\
2 & \mapsto & 013\\
3 & \mapsto & 24\\
4 & \mapsto & 035\\
5 & \mapsto & 46\\
 & \vdots & \\
 \end{array} \]
For this choice of $\mathbf{a}$, $C=2$ (longest run of 0s is 1), and hence $k=4$. It is clear that $| \varrho^4(a)|>2^4$ for any letter $a$.
Choose $F=\{ [0], [1], [2], [3] \}$.  Four ($k=4$) iterations of $\varrho$ applied to some letter (here: 5) successively yield
\[ 5 \mapsto\ 46 \mapsto \underline{0}35\underline{0}57 \mapsto \underline{0}\underline{1}
2446\underline{0}\underline{1}4668 \mapsto \underline{0}\underline{1}\underline{0}\underline{2}\underline{0}
13\underline{0}35\underline{0}57\underline{0}\underline{1}\underline{0}\underline{2}\underline{0}35\underline{0}57\underline{0}57\underline{0}79\]

\caption{An example of the coloring process in the proof of Theorem \ref{thm:primrecquasi}.
Red letters are underlined. After four iterations of $\varrho$ there remain only twelve unmarked letters
($12\leqslant 2^4$), thus at most twelve letters managed to escape $F$. In fact, only nine out of the twelve unmarked letters --- $5,5,7,5,7,5,7,7,9$ --- actually escaped $F$.
\label{tab:ex-red}}
\end{table}


Note that no red zero can create a letter outside of $F$ because we apply the substitution to red zeros at most $C+1$ times. Also, all red letters originate from red zeros so all letters outside of $F$ will be non-red in the end. But the number of non-red letters at most doubles because every letter is substituted with at most two non-zero letters. Thus, for every isolated point $[i]$, 
\[
\left|\left\{[j]\in\varrho^{C+2}_{\bf{a}}([i])\colon [j]\not\in F\right\} \right| \leqslant 2^{C+2}<2^{C+2}+1\leqslant r(M)^{C+2}, 
\]
and a similar inequality holds for all infinities in $\mathcal A$ due to the continuity of $\varrho_{\bf{a}}$. By Lemma \ref{prop:quasicompact}, this implies the quasicompactness of $M$ (after normalization).
%
%
%
%
%
%
%
%
%
%
%
%
%
\end{proof}
